%% ma: L11-B15-0003
%% lop: 11
%% chuong: 5
%% bai: 15
%% dang: 11.15.1
%% loai: TL
%% muc: VD
%% dap_an: "-"
%% nguon: "tai-lieu-hoc-tap-mon-toan-lop-11-theo-chuong-trinh-giao-duc-pho-thong-2018.pdf (Lớp Toán Thầy Duy, Phạm Lê Duy, THPT Chu Văn An), Chương 3 Bài 1, Dạng 1, Ví dụ 3 (tr.307)"
%% kien_thuc: "Lớp 11 Bài 15 (giới hạn của dãy số)"
%% trang_thai: cho-duyet
%% ly_do: "Dạng mới đề xuất (11.15.x), chờ giáo viên duyệt dạng; tài liệu gốc không có lời giải mẫu, đáp án do kho tự giải và kiểm bằng sympy."
\begin{ex}
Chứng minh rằng

a) $\lim\left(\dfrac{3\cdot3^n-\sin3n}{3^n}\right)=3$;

b) $\lim\left(\sqrt{n^2+n}-n\right)=\dfrac12$.
\loigiai{\sloppy a) $\dfrac{3\cdot3^n-\sin3n}{3^n}-3=\dfrac{-\sin3n}{3^n}$ và $\left|\dfrac{-\sin3n}{3^n}\right|\le\left(\dfrac13\right)^n\to0$, nên hiệu tiến về $0$.

b) $\sqrt{n^2+n}-n=\dfrac{n}{\sqrt{n^2+n}+n}=\dfrac{1}{\sqrt{1+\frac1n}+1}$, nên
\[\left(\sqrt{n^2+n}-n\right)-\dfrac12=\dfrac{1-\sqrt{1+\frac1n}}{2\left(\sqrt{1+\frac1n}+1\right)}\to\dfrac{1-1}{2\cdot2}=0.\]}
\end{ex}
